\subsection{Embedding e ricerca per similarità}\label{subsec:embedding}

Il primo problema che un sistema RAG deve risolvere è, dato un testo, come rappresentarlo in una forma che permetta
di confrontarne il significato con quello di una domanda. La soluzione standard è l'\emph{embedding}: una rete neurale
proietta ogni porzione di testo (un documento, o più spesso un suo frammento) in un vettore di numeri reali di
dimensione fissa, costruito in modo che frammenti dal significato simile vengano proiettati in punti vicini dello
spazio vettoriale. La vicinanza tra due vettori $\mathbf{u}$ e $\mathbf{v}$ si misura tipicamente con la
\emph{similarità coseno}, ovvero il coseno dell'angolo compreso tra i due vettori:
\begin{equation}
  \cos(\mathbf{u}, \mathbf{v}) = \frac{\mathbf{u} \cdot \mathbf{v}}{\lVert \mathbf{u} \rVert \, \lVert \mathbf{v} \rVert},
  \label{eq:coseno}
\end{equation}
che vale 1 per vettori con la stessa direzione e diminuisce man mano che le loro direzioni divergono.

Il punto di partenza per costruire modelli di embedding è in genere Bidirectional Encoder Representations
from Transformers (\emph{BERT}), proposto da Devlin et al.\ \cite{devlin2019bert}. A differenza dei LLM generativi descritti in precedenza, BERT utilizza solo la parte \emph{encoder} del Transformer: ogni token osserva
l'intera sequenza in entrambe le direzioni, e il modello produce per ciascun token un vettore che ne rappresenta il
significato nel contesto, senza generare testo. BERT viene pre-addestrato su grandi quantità di testo a ricostruire
parole mascherate casualmente all'interno delle frasi, e poi specializzato tramite \emph{fine-tuning} su compiti di
comprensione del linguaggio: la classificazione di un testo (per esempio l'analisi del sentiment), il riconoscimento
di entità nominate come persone e luoghi, l'individuazione della risposta a una domanda all'interno di un passaggio
dato, e la classificazione di coppie di frasi, per esempio per stabilire se due frasi sono parafrasi l'una dell'altra.

Nessuno di questi compiti, tuttavia, richiede di produrre un singolo vettore che rappresenti un'intera frase, e BERT
non è addestrato a farlo. Per confrontare due testi, BERT li riceve concatenati in un'unica sequenza e restituisce un
punteggio per la coppia, secondo un'architettura detta \emph{cross-encoder}. Reimers e Gurevych \cite{reimers2019sbert}
osservano che questo approccio è impraticabile per il recupero: ogni coppia richiede un passaggio completo attraverso
la rete, e trovare la coppia più simile in una collezione di 10.000 frasi richiede circa 50 milioni di passaggi, pari a
circa 65 ore di calcolo. D'altra parte, ricavare un vettore per ciascuna frase da BERT così com'è, per esempio facendo
la media dei vettori dei suoi token, produce embedding di bassa qualità.

La loro soluzione, Sentence-BERT (SBERT), affina BERT proprio per produrre embedding di frase confrontabili tramite la
similarità coseno. Durante l'addestramento, i due testi di ogni coppia non vengono concatenati, ma attraversano
separatamente la stessa rete, con gli stessi pesi; i due vettori risultanti vengono confrontati, e l'errore commesso
nel valutare la relazione tra i due testi (per esempio, se si tratta di parafrasi) viene usato per aggiornare i pesi.
Questa configurazione è detta \emph{siamese}, perché equivale a due copie gemelle della stessa rete. La condivisione dei
pesi garantisce che entrambi i testi vengano trasformati dalla stessa funzione, e quindi proiettati nello stesso
spazio, mentre l'addestramento fa sì che in questo spazio la similarità coseno rifletta la somiglianza di significato.
Una volta addestrata, la rete può calcolare l'embedding di ogni frammento una sola volta e indipendentemente dagli altri,
e il confronto con una domanda si riduce a una similarità coseno tra vettori già pronti: il tempo necessario a trovare
la coppia più simile tra 10.000 frasi scende così da circa 65 ore a circa 5 secondi.

Sentence-BERT è pensato per misurare la somiglianza tra frasi; nel recupero, invece, occorre confrontare una domanda
breve con passaggi di testo più lunghi che ne contengono la risposta. Karpukhin et al.\ \cite{karpukhin2020dpr}
applicano lo stesso principio a questo compito con il \emph{Dense Passage Retriever} (DPR), composto da due encoder
basati su BERT, uno per le domande e uno per i passaggi, addestrati su coppie domanda-passaggio in modo che l'embedding
di una domanda risulti vicino a quello del passaggio che contiene la risposta. Gli autori mostrano come DPR superi nettamente i metodi lessicali classici come BM25, che assegna un punteggio a ciascun documento in base alle parole che ha in
comune con la domanda, pesate secondo la loro frequenza nel documento e la loro rarità nella collezione.

Il vantaggio deriva dal fatto che un retriever basato su embedding confronta significati e non parole: una domanda come
``quando si paga la seconda rata?'' può essere associata a un passaggio che parla di ``scadenza del secondo versamento'',
pur avendo con esso quasi nessuna parola in comune. Per questo i retriever basati su embedding vengono detti
\emph{densi}, in contrapposizione alle rappresentazioni \emph{sparse} dei metodi lessicali come BM25, che descrivono un
testo con un vettore grande quanto l'intero vocabolario e composto quasi interamente da zeri. I metodi lessicali
restano comunque competitivi quando conta la corrispondenza esatta di un termine, come nel caso di codici o nomi propri.

Una domanda e un documento vengono quindi confrontati non piu sulla
base delle parole esatte che condividono, ma sulla distanza tra i rispettivi vettori nello spazio
degli embedding, calcolato in modo efficiente anche su milioni di candidati tramite strutture
dati specializzate (indici vettoriali) mantenute da un database vettoriale dedicato.

Un confronto esaustivo di un vettore query con ogni singolo vettore indicizzato (ricerca
esatta) diventa proibitivo già a poche centinaia di migliaia di frammenti, ben al di sotto della  scala di molti corpora documentali reali. I database vettoriali si affidano a una ricerca approssimata del vicino piu prossimo, che rinuncia alla garanzia di trovare sempre il risultato esatto in cambio di un tempo di
ricerca molto piu basso. Malkov e Yashunin \cite{malkov2018hnsw} propongono uno degli schemi più diffusi a
questo scopo, Hierarchical Navigable Small World (HNSW): un grafo multilivello in cui ogni
nodo rappresenta un vettore indicizzato, organizzato in strati via via piu densi, in cui la ricerca parte dallo strato piu rado (pochi nodi, collegamenti a lungo raggio) e scende progressivamente verso strati piu fitti, restringendo il raggio di ricerca a ogni passo; un meccanismo che ottiene un tempo di ricerca che cresce in modo logaritmico con il numero $N$ di vettori
indicizzati, $O(\log N)$, invece che linearmente, $O(N)$, come nella ricerca esatta.
